\subsection{A guide to the argument}\label{ms:guide}

The proof has two distinct jobs. It must show that a law with a large
Berry--Esseen ratio has nearly two-point geometry, and then compare its
sum with the corresponding binomial sum accurately enough to detect even
a very small improvement. The first job uses contact and Fourier
analysis; the second uses the discrete count structure.

\paragraph{Start at the first failure of the strict bound.}
Assume a sharp base through $41$, and choose the least $n\ge42$ with
$C_n\ge\CE$. Then $C_{n-1},C_{n-2}\le\CE\le C_n$, so
$d_n=n(C_{n-1}-C_n)_+=0$. The uniform smoothing bound is
\[
 C_n\le\sqrt{\frac n{n-1}}\,C_{n-1}
       +\frac{\sqrt n}{15(n-1)^{3/2}}
       +\frac{2\sqrt n}{25(n-1)^2}.
\]
It gives $0\le n(C_n-\CE)<3003/10250<0.3$.
The additive remainder in this normalized inequality is of order $1/n$.
Thus the current ratio is close to $\CE$, and both deleted indices have
sharp bounds. The floor $42$ is the one used by the proposed local
comparison; its choice fixes the finite base required by this route.

\paragraph{Use extremality to restrict the support.}
The supremum at a fixed $n$ is attained by a law $P$ and a closed threshold
$H$. Adding a small atom at $y$, then standardizing, gives the
first-variation inequality
\[
 \Pp(S_{n-1}\le H-y)\le Q_H(y),\qquad
 \Pp(S_{n-1}\le H-y)=Q_H(y)\quad(y\in\supp P),
\]
where $Q_H$ is the explicit piecewise cubic function in
\eqref{ms:contact-majorant}. Its cubic growth first gives compact support.
The quantitative step estimates the Gaussian part of contact at a radius
$R$, then proves that the upper bound decreases along both exterior rays.
Since the predecessor term vanishes, this yields the three support bounds
of Section~\ref{ms:support}, provided their moment bounds have been proved
in the order in which they are used.

\paragraph{Measure the departure from the extremal two-point law.}
Let $p=\Pp(X>0)$, $q=1-p$, and $m=\E X_+=\E X_-$.
Replacing each sign class by its conditional mean leaves a centered
conditional fluctuation. In suitable units this is
\[
 X=\ell(B_0+U-p),\qquad B_0=\ind_{\{X>0\}},\qquad
 \E(U\mid B_0)=0,\qquad \ell=m/(pq).
\]
The physical fluctuation is $W=\ell U$; the count-scale noise is
$s=\E U^2$, and $\lambda=ns$ is its variance over $n$ draws.

The zero-bias density is $\tau(t)=\E[X\ind_{\{X\ge t\}}]$.
It is constant on each gap $(a,b)$ in the support. Consequently
\[
 J=\sum_{(a,b)}\tau_{a,b}(b-a)^2
\]
is the expected length of the gap containing a zero-biased observation,
with gap length zero on the support. For a two-point law this is just
its span $\ell$, and for the law of~\eqref{ms:bernoulli-limit} the quantity
$\phi(0)(\gamma+3\ell)/(6\rho)$ is exactly the limit displayed there. The
inequality
\[
 D=(3+\sqrt{10})\rho-\gamma-3J\ge0
\]
says that $\phi(0)(\gamma+3J)/6$, which extends this first-order lattice
term to arbitrary support gaps, never exceeds $\CE\rho$.
Its equality case is exactly Esseen's oriented two-point law
(Corollary~\ref{ms:zero-defect}). More usefully, its nonnegative
 decomposition bounds the sign probability, inward endpoint movement
and conditional variance simultaneously.

\paragraph{Connect the defect to the error of normal approximation.}
Integrate the contact slack against $\tau$. Evaluating this integral
through support gaps and through the zero-bias transform of the sum
gives the conserved identity of Section~\ref{ms:conserved}. The identity
contains the score excess $n(C_n-\CE)\rho$, the defect cost and a
nonnegative leakage term. A Fourier remainder is the quantity to be
bounded. Once that remainder and the necessary sign conditions have
been controlled, one obtains a resource inequality of the form
\[
 n(C_n-\CE)\rho+\kappa nD+n^{3/2}L\le\mathcal B(H).
\]
Together with the defect decomposition, this inequality bounds the noise
and the inward displacement.
Retaining the first term, instead of discarding it, sharpens the budget.

\paragraph{Compare sums with the same Bernoulli skeleton.}
Conditioning on $K=\sum B_{0,i}\sim\operatorname{Bin}(n,p)$ gives
\[
 \Pp(S_n\le H)=\sum_{j=0}^n\binom njp^jq^{n-j}
                   \Pp(V_j\le k+u-j),
 \qquad np+H/\ell=k+u,
\]
where $V_j$ is the sum of $j$ independent copies of $U\mid B_0=1$
and $n-j$ independent copies of $U\mid B_0=0$. Subtracting
the binomial distribution function gives the exact score identity
\eqref{ms:master-identity}. At $\lambda=0$ the original law is two-point,
so Schulz's theorem applies directly. At positive noise, the outgoing
and incoming probabilities must be estimated with errors proportional
to $\lambda$, uniformly as $\lambda$ decreases to zero. The exact identity
makes this comparison possible; the required relative bounds are separate
steps in the argument. The three-candidate curvature formula simplifies
one of them.

\begin{figure}[htbp]
\centering
\begin{tikzpicture}[
 node distance=8mm,
 box/.style={draw,rounded corners=2pt,align=center,font=\small,
             text width=66mm,inner sep=5pt},
 input/.style={box,text width=48mm,dashed},
 arr/.style={-{Stealth[length=2mm]},semithick}]
\node[box] (entry) {First index with $C_n\ge\CE$, $n\ge42$\\
  sharp predecessors and zero loss};
\node[box,below=of entry] (contact) {Attainment and contact\\
  initial support bound};
\node[box,below=of contact] (resource) {Conserved identity\\
  joint defect and phase resources};
\node[box,below=of resource] (count) {Confinement and conditional counts\\
  comparison with the Bernoulli law};
\node[box,below=of count] (end) {Local exclusion and 13-role closure\\
  completes the global proof ($C_n\le\CE$, $C^*=\CE$)};
\node[input,right=12mm of entry] (base) {Verified finite base\\
  $C_j\le\CE$, $1\le j\le48$};
\node[input,right=12mm of resource] (fourier) {Valid Fourier enclosures\\
  and ordered moment bounds};
\node[input,right=12mm of count] (local) {Complete signed local cover\\
  and verified tail suite};
\draw[arr] (entry)--(contact);
\draw[arr] (contact)--(resource);
\draw[arr] (resource)--(count);
\draw[arr] (count)--(end);
\draw[arr,dashed] (base)--(entry);
\draw[arr,dashed] (fourier)--(resource);
\draw[arr,dashed] (local)--(count);
\end{tikzpicture}
\caption{The route to the local exclusion and global closure. Dashed boxes show
the certificate inputs; their completed \texttt{full\_pass} verification is
recorded in Section~\ref{ms:certificates}.}
\label{ms:fig:dependencies}
\end{figure}

\subsection{Notation and reading order}\label{ms:notation}
The main variables are listed below. Auxiliary variables used only in a
single estimate are defined where they occur.
\begin{center}\small
\begin{tabular}{@{}lP{110mm}@{}}\toprule
Symbol & Meaning\\\midrule
$H$, $z=H/\sqrt n$ & Raw and standardized thresholds.\\
$p,q,m$ & Sign probabilities and common one-sided first moment.\\
$\rho,\gamma,\beta$ & $\E|X|^3$, $\E X^3$, and $\E[X|X|]$.\\
$\ell$ & Span between conditional means; $\ell=m/(pq)$.\\
$W$, $U=W/\ell$ & Physical and count-scale conditional noise.\\
$\theta$, $s$, $\lambda$ & $\E W^2$, $\E U^2$, and $ns$; $\theta=s/(pq+s)$.\\
$D$, $d_n$, $d_Y$ & Moment defect, predecessor score decrement, and inward count-gap displacement.\\
$k,u$ & Nearest integer count and phase: $np+H/\ell=k+u$, $|u|\le1/2$.\\
$\sigma=\sqrt{10}$ & Constant in the defect decomposition.\\\bottomrule
\end{tabular}
\end{center}

Sections~\ref{ms:extremizers}--\ref{ms:counts} contain the analytic
argument, and Section~\ref{ms:certificates} the main global theorem
and the \texttt{full\_pass} verification of the thirteen certificate roles. The appendices contain the
longer Fourier estimates, the finite-index arguments and the explicit
comparison formulas.
